\documentclass[10pt]{article}
\usepackage[spanish]{babel}
\usepackage[utf8]{inputenc}
\usepackage[T1]{fontenc}
\usepackage[margin=2cm]{geometry}
\usepackage{enumitem}
\title{Propuesta de proyecto: Sombra digital de tráfico en una intersección de Paute}
\author{David Timbi y equipo}
\date{\today}
\begin{document}\maketitle
\section*{Problema y motivación}
La gestión local requiere observar cómo cambia el flujo vehicular en intersecciones concurridas. Un modelo microscópico aislado puede reproducir circulación plausible, pero no representa por sí solo las condiciones observadas en calle. Se propone conectar una cámara de bajo costo con una escena de simulación de Paute y hacer visible la correspondencia entre detección y entidad digital.
\section*{Objetivo}
Construir y evaluar un prototipo de sombra digital para Av. Interoceánica (E40) y José Miguel Barzallo, junto al Mercado 26 de Febrero, que detecte/clasifique vehículos, estime su sentido y los refleje en una visualización 3D apoyada en SUMO.
\section*{Método y alcance}
El video se procesa en el equipo local con YOLOv11 y seguimiento; eventos direccionales se transmiten por WebSocket al cliente Three.js. La red vial y edificios se derivan de OpenStreetMap y SUMO. Se reportarán métricas solo después de crear una referencia humana independiente: emparejamiento uno a uno por objeto y ventana temporal, precision/recall, clase, sentido, demora y error espacial. SUMO sintético será línea base/soporte de software, claramente separado de observación real.
\section*{Entregables y validación}
Aplicación local (HTTP 8085, WebSocket de telemetría), red 3D exportable, escenario SUMO, código y documentación reproducible, informe IEEE y defensa. La validación de campo queda condicionada a anotar video consentido, identificar intervalos, congelar versiones y publicar los denominadores. La versión actual es un prototipo; todavía no hay conjunto anotado ni resultados estadísticos de campo.
\end{document}
